\SourceLeaf{240}%
\[
\begin{aligned}
aa+bb+cc&=E\\
aa'+bb'+cc'&=F\\
a'a'+b'b'+c'c'&=G\\
a\alpha+b\beta+c\gamma&=m\qquad .\ .\ .\ .\ .\ .\ .\ .\ .\ .\ .\quad(4)\\
a\alpha'+b\beta'+c\gamma'&=m'\qquad .\ .\ .\ .\ .\ .\ .\ .\ .\ .\ .\quad(5)\\
a\alpha''+b\beta''+c\gamma''&=m''\qquad .\ .\ .\ .\ .\ .\ .\ .\ .\ .\ .\quad(6)\\
a'\alpha+b'\beta+c'\gamma&=n\qquad .\ .\ .\ .\ .\ .\ .\ .\ .\ .\ .\quad(7)\\
a'\alpha'+b'\beta'+c'\gamma'&=n'\qquad .\ .\ .\ .\ .\ .\ .\ .\ .\ .\ .\ .\quad(8)\\
a'\alpha''+b'\beta''+c'\gamma''&=n''\qquad .\ .\ .\ .\ .\ .\ .\ .\ .\ .\ .\quad(9)\\
AA+BB+CC&=EG-FF=\Delta
\end{aligned}
\]
\par\indent Eliminieren wir aus den Gleichungen 1, 4, 7 die Größen \(\beta,\gamma\), indem wir die Gleichungen der Reihe nach mit \(bc'-cb',\ b'C-c'B,\ cB-bC\) multiplizieren und addieren. So erhalten wir
\[
\begin{aligned}
&(A(bc'-cb')+a(b'C-c'B)+a'(cB-bC))\alpha\\
&\quad=D(bc'-cb')+m(b'C-c'B)+n(cB-bC)
\end{aligned}
\]
eine Gleichung, die wir leicht in die folgende Form bringen:
\[
AD=\alpha\Delta+a(nF-mG)+a'(mF-nE)
\]
Entsprechend ergibt die Elimination der Größen \(\alpha,\gamma\) oder \(\alpha,\beta\) aus denselben Gleichungen
\[
\begin{aligned}
BD&=\beta\Delta+b(nF-mG)+b'(mF-nE)\\
CD&=\gamma\Delta+c(nF-mG)+c'(mF-nE)
\end{aligned}
\]
Multiplizieren wir diese drei Gleichungen mit \(\alpha'',\beta'',\gamma''\) und addieren sie, so erhalten wir
\[
\begin{aligned}
DD''&=(\alpha\alpha''+\beta\beta''+\gamma\gamma'')\Delta\\
&\quad+m''(nF-mG)+n''(mF-nE)\quad .\ .\ .\ .\quad(10)
\end{aligned}
\]
\par\indent Verfahren wir ebenso mit den Gleichungen 2, 5, 8, so ergibt sich
\[
\begin{aligned}
AD'&=\alpha'\Delta+a(n'F-m'G)+a'(m'F-n'E)\\
BD'&=\beta'\Delta+b(n'F-m'G)+b'(m'F-n'E)\\
CD'&=\gamma'\Delta+c(n'F-m'G)+c'(m'F-n'E)
\end{aligned}
\]
Multipliziert man diese Gleichungen mit \(\alpha',\beta',\gamma'\) und addiert sie, so erhält man:
\[
\begin{aligned}
D'D'&=(\alpha'\alpha'+\beta'\beta'+\gamma'\gamma')\Delta\\
&\quad+m'(n'F-m'G)+n'(m'F-n'E)
\end{aligned}
\]
